% Illustrative Application

We demonstrated signal hierarchy theory's applicability through \textit{Bacillus subtilis} sporulation~\cite{subtilis2026}, illustrating quantitative threshold determination inexpressible in classical Bio-PN. Full model details, parameter derivations, and comprehensive sensitivity analysis are available in~\cite{subtilis2026}.

\subsection{\textit{Bacillus subtilis} Sporulation: ATP-Gated Developmental Commitment}

\vspace{3pt}
\noindent Sporulation represents a canonical energy-controlled developmental decision where metabolic state (ATP concentration) hierarchically regulates gene expression networks. Under nutrient starvation, cells must commit to either sporulation (energy-intensive but stress-resistant) or competence (DNA uptake for genetic adaptation). This commitment is irreversible once initiated, requiring formal decision finality semantics.

\vspace{3pt}
\noindent\textbf{Hierarchical Architecture}:

\vspace{3pt}
\noindent\textbf{Layer 3 (Metabolic State)}: ATP production via glycolysis and oxidative phosphorylation. ATP functions simultaneously as metabolic product (normal arc output) and regulatory signal (signal flow arc source at higher layers).

\vspace{3pt}
\noindent\textbf{Layer 2 (Energy Sensing)}: ATP threshold gates phosphorelay cascade activation. Signal consumption semantics: ATP depletion below commitment threshold irreversibly blocks sporulation pathway.

\vspace{3pt}
\noindent\textbf{Layer 1 (Phosphorelay Integration)}: KinA/KinB $\to$ Spo0F $\to$ Spo0A cascade integrates multiple stress inputs when ATP sufficient. Spo0A-P accumulation triggers layer 0 commitment.

\vspace{3pt}
\noindent\textbf{Layer 0 (Execution Machinery)}: $\sigma^F$/$\sigma^E$ sigma factors implement sporulation vs. competence pathways. Exclusive commitment enforced by signal consumption at layer 2.

\subsection{Signal Hierarchy Model}

\vspace{3pt}
\noindent\textbf{Signal places ($\Psi$)}: ATP (layer 3), phosphorelay activation signal (layer 2), Spo0A-P commitment signal (layer 1), sigma factor availability (layer 0).

\vspace{3pt}
\noindent\textbf{Signal flow arcs}: ATP $\to$ phosphorelay activation (consumption semantics), Spo0A-P $\to$ sigma factor deployment (consumption semantics).

\vspace{3pt}
\noindent\textbf{Enablement thresholds}: Sporulation requires $E(\text{phosphorelay}, \text{ATP}) = 2.5$ mM. ATP below threshold activates competence.

\vspace{3pt}
\noindent\textbf{Contrast with test arc model}: Classical Bio-PN uses test arcs for ATP regulation (non-consuming reads). This allows sporulation to activate even at low ATP, contradicting biological reality where energy insufficiency irreversibly blocks the pathway. Test arcs cannot express decision finality---the system remains reversible at all ATP concentrations.

\subsection{Quantitative Threshold Determination}

\vspace{3pt}
\noindent\textbf{Analysis method}: Basin boundary computation via reachability analysis. Starting from nutrient-depleted initial state, systematically vary ATP concentration and compute attractor convergence. Commitment threshold identifies ATP level where basin geometry transitions from competence dominance to sporulation dominance.

\begin{figure}[t]
\centering
\includegraphics[width=0.95\columnwidth]{figures/bacillus_atp_threshold.pdf}
\caption{\textbf{ATP commitment threshold in \textit{B. subtilis} sporulation.} Signal hierarchy analysis computes developmental decision threshold through basin boundary determination. The theoretical framework predicts commitment at 2.38 mM ATP (blue dashed line), matching experimental measurement of 2.21 mM (red dashed line, \cite{fujita2005}) with 7\% error. Below threshold, competence pathway dominates (94\% reliability); above threshold, sporulation pathway dominates (91\% reliability). Classical Bio-PN with test arcs cannot compute this threshold---non-consuming reads provide no consumption semantics for basin boundary analysis. Curve shows sigmoid fit to 1000 stochastic trajectories; shaded regions indicate pathway dominance zones. This quantitative prediction demonstrates signal hierarchy's capability for computing decision thresholds inexpressible in classical Petri net theory.}
\label{fig:atp_threshold}
\end{figure}

\vspace{3pt}
\noindent\textbf{Theoretical prediction}: The theoretical framework predicts ATP commitment threshold at $2.38 \pm 0.15$ mM (mean $\pm$ SD over 1000 stochastic trajectories, Figure~\ref{fig:atp_threshold}). Below this threshold, competence pathway achieves 94\% reliability. Above threshold, sporulation pathway achieves 91\% reliability.

\vspace{3pt}
\noindent\textbf{Experimental validation}: Fujita \& Losick~\cite{fujita2005} measured sporulation commitment threshold in glucose-limited chemostats at $2.21 \pm 0.18$ mM ATP. Our model achieves 7\% prediction error (Figure~\ref{fig:atp_threshold}).

\vspace{3pt}
\noindent\textbf{Classical Bio-PN limitation}: Test arc models cannot compute this threshold. Non-consuming reads provide no consumption semantics, preventing basin boundary determination. The question ``what ATP level triggers irreversible commitment?'' is formally inexpressible in classical Bio-PN theory.

\subsection{Biological Significance}

\vspace{3pt}
\noindent This validated example demonstrates signal hierarchy's predictive capability for energy-controlled developmental decisions. ATP functions simultaneously as:

\vspace{3pt}
\noindent\textbf{Metabolite}: Produced by glycolysis (normal arc output), consumed by biosynthesis (normal arc input)

\vspace{3pt}
\noindent\textbf{Regulatory signal}: Gates phosphorelay activation (signal flow arc with consumption), hierarchically controlling developmental pathway accessibility

\vspace{3pt}
\noindent Classical Bio-PN cannot formally express this dual role. Signal hierarchy provides the first framework enabling quantitative prediction of commitment thresholds through integrated metabolic-regulatory modeling with consumption semantics.

\vspace{3pt}
\noindent\textbf{Generalizability}: The formalism applies to any energy-gated developmental decision: \textit{C. elegans} dauer formation (energy status gates dauer pathway), yeast pseudohyphal growth (glucose availability controls filamentous switching), bacterial persister formation (ATP depletion triggers dormancy). All exhibit hierarchical organization where metabolic state preempts regulatory networks through signal consumption.
